\chapter{第2章：Eisenstein级数与解析工具 习题}
\difficultykey

\section{基础题 \easy}

\begin{problem}{Bernoulli数表}
\easy
由生成函数
\[
\frac{t}{e^t-1}=\sum_{n=0}^{\infty} B_n\frac{t^n}{n!}
\]
求出 $B_0,B_1,B_2,B_3,B_4,B_5,B_6$。
\end{problem}
\begin{solution}
把
\[
e^t-1=t+\frac{t^2}{2}+\frac{t^3}{6}+\frac{t^4}{24}+\cdots
\]
代入并做幂级数除法，可得
\[
\frac{t}{e^t-1}=1-\frac{t}{2}+\frac{t^2}{12}-\frac{t^4}{720}+\frac{t^6}{30240}+\cdots.
\]
与
\[
\sum_{n=0}^{\infty} B_n\frac{t^n}{n!}
= B_0+B_1 t+\frac{B_2}{2!}t^2+\frac{B_3}{3!}t^3+\frac{B_4}{4!}t^4+\frac{B_5}{5!}t^5+\frac{B_6}{6!}t^6+\cdots
\]
逐项比较，得到
\[
B_0=1,\qquad B_1=-\frac12,\qquad B_2=\frac16,
\]
\[
B_3=0,\qquad B_4=-\frac1{30},\qquad B_5=0,\qquad B_6=\frac1{42}.
\]
\end{solution}

\begin{problem}{偶整数处的zeta值}
\easy
利用公式
\[
\zeta(2k)=\frac{(2\pi)^{2k}|B_{2k}|}{2(2k)!}
\]
计算 $\zeta(2)$、$\zeta(4)$ 与 $\zeta(6)$。
\end{problem}
\begin{solution}
当 $k=1$ 时，$B_2=1/6$，所以
\[
\zeta(2)=\frac{(2\pi)^2\cdot \frac16}{2\cdot 2!}
=\frac{4\pi^2}{24}=\frac{\pi^2}{6}.
\]

当 $k=2$ 时，$|B_4|=1/30$，故
\[
\zeta(4)=\frac{(2\pi)^4\cdot \frac1{30}}{2\cdot 4!}
=\frac{16\pi^4}{30\cdot 48}=\frac{\pi^4}{90}.
\]

当 $k=3$ 时，$B_6=1/42$，故
\[
\zeta(6)=\frac{(2\pi)^6\cdot \frac1{42}}{2\cdot 6!}
=\frac{64\pi^6}{42\cdot 1440}=\frac{\pi^6}{945}.
\]
所以
\[
\zeta(2)=\frac{\pi^2}{6},\qquad
\zeta(4)=\frac{\pi^4}{90},\qquad
\zeta(6)=\frac{\pi^6}{945}.
\]
\end{solution}

\begin{problem}{除数函数计算}
\easy
计算下列值：
\[
\sigma_1(12),\qquad \sigma_3(12),\qquad \sigma_5(6).
\]
\end{problem}
\begin{solution}
$12$ 的正因子为 $1,2,3,4,6,12$，故
\[
\sigma_1(12)=1+2+3+4+6+12=28.
\]
又
\[
\sigma_3(12)=1^3+2^3+3^3+4^3+6^3+12^3
=1+8+27+64+216+1728=2044.
\]

$6$ 的正因子为 $1,2,3,6$，所以
\[
\sigma_5(6)=1^5+2^5+3^5+6^5=1+32+243+7776=8052.
\]
\end{solution}

\begin{problem}{乘性的例子}
\easy
利用除数函数的乘性，计算
\[
\sigma_1(15),\qquad \sigma_3(35).
\]
\end{problem}
\begin{solution}
因为 $15=3\cdot 5$ 且 $(3,5)=1$，所以
\[
\sigma_1(15)=\sigma_1(3)\sigma_1(5)=(1+3)(1+5)=4\cdot 6=24.
\]

又因为 $35=5\cdot 7$ 且 $(5,7)=1$，所以
\[
\sigma_3(35)=\sigma_3(5)\sigma_3(7)=(1+5^3)(1+7^3)=(1+125)(1+343).
\]
故
\[
\sigma_3(35)=126\cdot 344=43344.
\]
\end{solution}

\begin{problem}{E4展开}
\easy
写出归一化 Eisenstein 级数 $E_4$ 的前四个非零项。
\end{problem}
\begin{solution}
由一般公式
\[
E_k(\tau)=1-\frac{2k}{B_k}\sum_{n\ge 1}\sigma_{k-1}(n)q^n,
\qquad q=e^{2\pi i\tau},
\]
对 $k=4$ 有 $B_4=-1/30$，所以
\[
E_4(\tau)=1+240\sum_{n\ge1}\sigma_3(n)q^n.
\]
又
\[
\sigma_3(1)=1,\quad \sigma_3(2)=9,\quad \sigma_3(3)=28,\quad \sigma_3(4)=73.
\]
因此
\[
E_4(\tau)=1+240q+2160q^2+6720q^3+17520q^4+\cdots.
\]
\end{solution}

\begin{problem}{E6与E8展开}
\easy
写出 $E_6$ 与 $E_8$ 的前四个非零项。
\end{problem}
\begin{solution}
对 $E_6$，由于 $B_6=1/42$，所以
\[
E_6(\tau)=1-504\sum_{n\ge1}\sigma_5(n)q^n.
\]
又
\[
\sigma_5(1)=1,\quad \sigma_5(2)=33,\quad \sigma_5(3)=244,\quad \sigma_5(4)=1057.
\]
故
\[
E_6(\tau)=1-504q-16632q^2-122976q^3-532728q^4+\cdots.
\]

对 $E_8$，由于 $B_8=-1/30$，所以
\[
E_8(\tau)=1+480\sum_{n\ge1}\sigma_7(n)q^n.
\]
又
\[
\sigma_7(1)=1,\quad \sigma_7(2)=129,\quad \sigma_7(3)=2188,\quad \sigma_7(4)=16513.
\]
故
\[
E_8(\tau)=1+480q+61920q^2+1050240q^3+7926240q^4+\cdots.
\]
\end{solution}

\begin{problem}{平移不变性}
\easy
利用 $q=e^{2\pi i\tau}$ 证明 Eisenstein 级数满足
\[
E_k(\tau+1)=E_k(\tau).
\]
\end{problem}
\begin{solution}
因为
\[
q(\tau+1)=e^{2\pi i(\tau+1)}=e^{2\pi i\tau}e^{2\pi i}=q(\tau),
\]
所以 $q$ 在变换 $\tau\mapsto \tau+1$ 下不变。

而 $E_k$ 的 $q$-展开形如
\[
E_k(\tau)=1+\sum_{n\ge1} a_n q^n,
\]
因此
\[
E_k(\tau+1)=1+\sum_{n\ge1} a_n q(\tau+1)^n
=1+\sum_{n\ge1} a_n q(\tau)^n
=E_k(\tau).
\]
这就证明了平移不变性。
\end{solution}

\begin{problem}{theta函数的对称性}
\easy
设
\[
\theta(t)=\sum_{n\in\Z} e^{-\pi n^2 t},\qquad t>0.
\]
利用 Poisson 求和公式证明
\[
\theta(t)=t^{-1/2}\theta(1/t),
\]
并据此写出 $\theta(4)$ 与 $\theta(1/4)$ 的关系。
\end{problem}
\begin{solution}
对高斯函数应用 Poisson 求和公式，可得经典恒等式
\[
\theta(t)=t^{-1/2}\theta(1/t).
\]
这是 theta 函数最基本的变换公式。

把 $t=4$ 代入，得到
\[
\theta(4)=4^{-1/2}\theta(1/4)=\frac12\theta(1/4).
\]
也就是说
\[
\theta(1/4)=2\theta(4).
\]
\end{solution}

\begin{problem}{模四特征}
\easy
定义模 $4$ 的 Dirichlet 特征
\[
\chi_4(n)=
\begin{cases}
0,& n\equiv 0,2\pmod 4,\\
1,& n\equiv 1\pmod 4,\\
-1,& n\equiv 3\pmod 4.
\end{cases}
\]
写出 $L(s,\chi_4)$ 的前几项，并说明哪些整数项消失。
\end{problem}
\begin{solution}
按定义
\[
L(s,\chi_4)=\sum_{n\ge1}\frac{\chi_4(n)}{n^s}.
\]
因为偶数上 $\chi_4(n)=0$，所以所有偶数项都消失；奇数上符号按模 $4$ 的余数交替。于是
\[
L(s,\chi_4)=1-\frac1{3^s}+\frac1{5^s}-\frac1{7^s}+\frac1{9^s}-\frac1{11^s}+\cdots.
\]
消失的正是所有偶数项。
\end{solution}

\section{进阶题 \medium}

\begin{problem}{除数函数乘性证明}
\medium
证明：若 $(m,n)=1$，则
\[
\sigma_r(mn)=\sigma_r(m)\sigma_r(n).
\]
\end{problem}
\begin{solution}
因为 $(m,n)=1$，$mn$ 的每个正因子都可以唯一写成 $d_1d_2$ 的形式，其中 $d_1\mid m$，$d_2\mid n$。因此
\[
\sigma_r(mn)=\sum_{d\mid mn} d^r
=\sum_{d_1\mid m}\sum_{d_2\mid n}(d_1d_2)^r.
\]
把幂拆开得
\[
\sigma_r(mn)=\left(\sum_{d_1\mid m} d_1^r\right)
\left(\sum_{d_2\mid n} d_2^r\right)
=\sigma_r(m)\sigma_r(n).
\]
这就证明了乘性。
\end{solution}

\begin{problem}{归一化系数公式}
\medium
证明对偶整数 $k\ge 4$，归一化 Eisenstein 级数满足
\[
E_k(\tau)=1-\frac{2k}{B_k}\sum_{n\ge1}\sigma_{k-1}(n)q^n.
\]
\end{problem}
\begin{solution}
章节中给出的公式是
\[
E_k(\tau)=1+\frac{2}{\zeta(1-k)}\sum_{n\ge1}\sigma_{k-1}(n)q^n.
\]
又对偶整数 $k\ge 2$ 有
\[
\zeta(1-k)=-\frac{B_k}{k}.
\]
因此
\[
\frac{2}{\zeta(1-k)}=\frac{2}{-B_k/k}=-\frac{2k}{B_k}.
\]
代回上式便得
\[
E_k(\tau)=1-\frac{2k}{B_k}\sum_{n\ge1}\sigma_{k-1}(n)q^n.
\]
这正是所求公式。
\end{solution}

\begin{problem}{E10与E12展开}
\medium
写出 $E_{10}$ 与 $E_{12}$ 的前两个非零 Fourier 系数。
\end{problem}
\begin{solution}
对 $k=10$，有 $B_{10}=5/66$，所以
\[
E_{10}(\tau)=1-264\sum_{n\ge1}\sigma_9(n)q^n.
\]
又
\[
\sigma_9(1)=1,\qquad \sigma_9(2)=1+2^9=513.
\]
因此
\[
E_{10}(\tau)=1-264q-135432q^2+\cdots.
\]

对 $k=12$，有 $B_{12}=-691/2730$，所以
\[
E_{12}(\tau)=1+\frac{65520}{691}\sum_{n\ge1}\sigma_{11}(n)q^n.
\]
又
\[
\sigma_{11}(1)=1,
\qquad
\sigma_{11}(2)=1+2^{11}=2049.
\]
因此
\[
E_{12}(\tau)=1+\frac{65520}{691}q+\frac{134250480}{691}q^2+\cdots.
\]
这两个例子说明：当权增长时，Fourier 系数依然由 Bernoulli 数与除数函数统一控制；真正新的现象并不在“是否有 $q$-展开”，而在不同权空间的维数和尖点部分如何出现。对 $k=10$，这里只有 Eisenstein 方向；而到 $k=12$，由于同时出现了 $\Delta$，模形式空间开始出现非平凡的 cusp 部分。
\end{solution}

\begin{problem}{权十二恒等式}
\medium
利用 $M_{12}$ 是二维空间，证明
\[
E_4^3-E_6^2=1728\Delta.
\]
\end{problem}
\begin{solution}
$E_4^3$ 与 $E_6^2$ 都是权 $12$ 模形式，而且它们的常数项都等于 $1$。因此
\[
E_4^3-E_6^2
\]
是权 $12$ 的尖点形式。

又因为 $M_{12}$ 是二维，而 $E_{12}$ 提供了一条 Eisenstein 方向，所以尖点空间 $S_{12}$ 是一维的。$\Delta$ 是一个非零的权 $12$ 尖点形式，因此
\[
S_{12}=\C\Delta.
\]
于是存在常数 $c$ 使
\[
E_4^3-E_6^2=c\Delta.
\]

接着比较 $q$ 的一次项。由前面已经算出的展开，
\[
E_4=1+240q+\cdots,
\qquad
E_6=1-504q+\cdots.
\]
所以
\[
E_4^3=1+720q+\cdots,
\qquad
E_6^2=1-1008q+\cdots.
\]
故
\[
E_4^3-E_6^2=1728q+\cdots.
\]
另一方面，$\Delta=q+O(q^2)$，所以必有 $c=1728$。因此
\[
E_4^3-E_6^2=1728\Delta.
\]
\end{solution}

\begin{problem}{Delta的首项系数}
\medium
由恒等式
\[
\Delta=\frac{E_4^3-E_6^2}{1728}
\]
计算 $\Delta$ 的前四个非零项。
\end{problem}
\begin{solution}
先写出
\[
E_4=1+240q+2160q^2+6720q^3+17520q^4+\cdots,
\]
\[
E_6=1-504q-16632q^2-122976q^3-532728q^4+\cdots.
\]
直接展开可得
\[
E_4^3=1+720q+179280q^2+16954560q^3+396974160q^4+\cdots,
\]
\[
E_6^2=1-1008q+220752q^2+16519104q^3+399517776q^4+\cdots.
\]
两者相减得到
\[
E_4^3-E_6^2=1728q-41472q^2+435456q^3-2543616q^4+\cdots.
\]
除以 $1728$ 便有
\[
\Delta=q-24q^2+252q^3-1472q^4+\cdots.
\]
\end{solution}

\begin{problem}{E2的变换公式}
\medium
已知
\[
E_2(\tau+1)=E_2(\tau),
\qquad
E_2\!\left(-\frac1\tau\right)=\tau^2E_2(\tau)+\frac{6}{\pi i}\tau.
\]
证明对任意
\[
\gamma=\begin{pmatrix}a&b\\c&d\end{pmatrix}\in\SL_2(\Z)
\]
都有
\[
E_2(\gamma\tau)=(c\tau+d)^2E_2(\tau)+\frac{6}{\pi i}c(c\tau+d).
\]
\end{problem}
\begin{solution}
定义
\[
F_\gamma(\tau)=E_2(\gamma\tau)-(c\tau+d)^2E_2(\tau)-\frac{6}{\pi i}c(c\tau+d).
\]
我们要证 $F_\gamma\equiv 0$。

对生成元 $T=\begin{psmallmatrix}1&1\\0&1\end{psmallmatrix}$，因为 $c=0,d=1$，有
\[
F_T(\tau)=E_2(\tau+1)-E_2(\tau)=0.
\]

对生成元 $S=\begin{psmallmatrix}0&-1\\1&0\end{psmallmatrix}$，公式正是已知条件，所以
\[
F_S(\tau)=0.
\]

现在设该公式已对 $\gamma_1,\gamma_2$ 成立。把 $\gamma_1\gamma_2$ 写成
\[
\gamma_1\gamma_2=
\begin{pmatrix}a&b\\c&d\end{pmatrix}.
\]
那么
\[
E_2(\gamma_1\gamma_2\tau)
= (c_1\gamma_2\tau+d_1)^2E_2(\gamma_2\tau)+\frac{6}{\pi i}c_1(c_1\gamma_2\tau+d_1).
\]
再对 $E_2(\gamma_2\tau)$ 代入同样公式并整理，恰好得到
\[
E_2(\gamma_1\gamma_2\tau)=(c\tau+d)^2E_2(\tau)+\frac{6}{\pi i}c(c\tau+d).
\]
也就是说，若公式对两个矩阵成立，则对它们的乘积也成立。

由于 $\SL_2(\Z)$ 由 $S,T$ 生成，而公式已对 $S,T$ 成立，所以对一切 $\gamma\in\SL_2(\Z)$ 都成立。
\end{solution}

\begin{problem}{E2的实解析修正}
\medium
定义
\[
E_2^*(\tau)=E_2(\tau)-\frac{3}{\pi\Im\tau}.
\]
证明 $E_2^*$ 满足真正的权 $2$ 模变换律。
\end{problem}
\begin{solution}
设 $\tau=x+iy$，并记
\[
\gamma=\begin{pmatrix}a&b\\c&d\end{pmatrix}.
\]
由上一题，
\[
E_2(\gamma\tau)=(c\tau+d)^2E_2(\tau)+\frac{6}{\pi i}c(c\tau+d).
\]
又由虚部公式，
\[
\Im(\gamma\tau)=\frac{y}{|c\tau+d|^2}.
\]
因此
\[
E_2^*(\gamma\tau)
=E_2(\gamma\tau)-\frac{3}{\pi\Im(\gamma\tau)}
=(c\tau+d)^2E_2(\tau)+\frac{6}{\pi i}c(c\tau+d)-\frac{3|c\tau+d|^2}{\pi y}.
\]
另一方面，
\[
(c\tau+d)^2-|c\tau+d|^2=(c\tau+d)(c\tau+d-c\bar\tau-d)=2iy\,c(c\tau+d).
\]
故
\[
-\frac{3|c\tau+d|^2}{\pi y}
=-\frac{3(c\tau+d)^2}{\pi y}+\frac{6i}{\pi}c(c\tau+d).
\]
注意到
\[
\frac{6}{\pi i}=-\frac{6i}{\pi},
\]
所以上式中的附加项正好抵消，得到
\[
E_2^*(\gamma\tau)=(c\tau+d)^2\left(E_2(\tau)-\frac{3}{\pi y}\right)
=(c\tau+d)^2E_2^*(\tau).
\]
这说明 $E_2^*$ 确实按权 $2$ 变换。
\end{solution}

\begin{problem}{Euler乘积}
\medium
证明：当 $\Re(s)>1$ 时，Dirichlet $L$-函数满足 Euler 乘积
\[
L(s,\chi)=\prod_p \left(1-\chi(p)p^{-s}\right)^{-1}.
\]
\end{problem}
\begin{solution}
由定义
\[
L(s,\chi)=\sum_{n\ge1}\frac{\chi(n)}{n^s}.
\]
在 $\Re(s)>1$ 时该级数绝对收敛。又因为 Dirichlet 特征是完全乘法的，所以对于每个素数 $p$，有几何级数展开
\[
\left(1-\chi(p)p^{-s}\right)^{-1}=1+\chi(p)p^{-s}+\chi(p)^2p^{-2s}+\cdots.
\]
把所有素数的这些级数相乘，依据算术基本定理，每个正整数
\[
n=\prod_p p^{\nu_p(n)}
\]
恰好出现一次，其对应项为
\[
\prod_p \chi(p)^{\nu_p(n)}p^{-\nu_p(n)s}=\chi(n)n^{-s}.
\]
因此乘积展开后得到
\[
\prod_p \left(1-\chi(p)p^{-s}\right)^{-1}
=\sum_{n\ge1}\frac{\chi(n)}{n^s}=L(s,\chi).
\]
\end{solution}

\begin{problem}{模四Gauss和}
\medium
对模 $4$ 的特征 $\chi_4$，计算它的 Gauss 和
\[
\tau(\chi_4)=\sum_{a=1}^{4}\chi_4(a)e^{2\pi i a/4}.
\]
\end{problem}
\begin{solution}
按定义逐项计算：
\[
\chi_4(1)=1,\qquad \chi_4(2)=0,\qquad \chi_4(3)=-1,\qquad \chi_4(4)=0.
\]
因此
\[
\tau(\chi_4)=e^{2\pi i/4}-e^{6\pi i/4}=i-(-i)=2i.
\]
故
\[
\tau(\chi_4)=2i,
\qquad
|\tau(\chi_4)|=2.
\]
\end{solution}

\section{挑战题 \hard}

\begin{problem}{zeta函数的函数方程}
\hard
设
\[
\Lambda(s)=\pi^{-s/2}\Gamma\!\left(\frac{s}{2}\right)\zeta(s).
\]
利用 theta 函数的变换公式证明
\[
\Lambda(s)=\Lambda(1-s).
\]
\end{problem}
\begin{solution}
先对 $\Re(s)>1$ 写 Mellin 变换：
\[
\Lambda(s)=\frac12\int_0^{\infty}(\theta(t)-1)t^{s/2-1}\,dt.
\]
把积分拆成两段：
\[
\Lambda(s)=\frac12\int_0^1(\theta(t)-1)t^{s/2-1}dt
+\frac12\int_1^{\infty}(\theta(t)-1)t^{s/2-1}dt.
\]
在第一段中令 $t=1/u$，并用
\[
\theta(1/u)=u^{1/2}\theta(u),
\]
得到
\[
\int_0^1(\theta(t)-1)t^{s/2-1}dt
=\int_1^{\infty}(u^{1/2}\theta(u)-1)u^{-s/2-1}du.
\]
把右边拆开并整理，可得
\[
\int_0^1(\theta(t)-1)t^{s/2-1}dt
=\int_1^{\infty}(\theta(u)-1)u^{(1-s)/2-1}du+\frac{2}{s-1}-\frac{2}{s}.
\]
于是
\[
\Lambda(s)=\frac{1}{s(s-1)}+\frac12\int_1^{\infty}(\theta(t)-1)
\left(t^{s/2-1}+t^{(1-s)/2-1}\right)dt.
\]
这个表达式对 $s$ 与 $1-s$ 完全对称，因此
\[
\Lambda(s)=\Lambda(1-s).
\]
这就是 Riemann zeta 函数的函数方程。
\end{solution}

\begin{problem}{L值与反正切}
\hard
证明
\[
L(1,\chi_4)=\frac{\pi}{4}.
\]
\end{problem}
\begin{solution}
由定义
\[
L(1,\chi_4)=1-\frac13+\frac15-\frac17+\cdots.
\]
另一方面，初等分析中的反正切级数给出
\[
\arctan x=x-\frac{x^3}{3}+\frac{x^5}{5}-\frac{x^7}{7}+\cdots,
\qquad |x|\le 1.
\]
令 $x=1$，便得到 Leibniz 级数
\[
\arctan 1=1-\frac13+\frac15-\frac17+\cdots.
\]
而 $\arctan 1=\pi/4$，故
\[
L(1,\chi_4)=\frac{\pi}{4}.
\]
\end{solution}

\begin{problem}{Gauss和的扭转性质}
\hard
设 $\chi$ 是模 $m$ 的 Dirichlet 特征，定义
\[
\tau(\chi)=\sum_{a\bmod m}\chi(a)e^{2\pi i a/m}.
\]
证明：若 $(n,m)=1$，则
\[
\sum_{a\bmod m}\chi(a)e^{2\pi i an/m}=\overline{\chi}(n)\tau(\chi).
\]
\end{problem}
\begin{solution}
由于 $(n,m)=1$，$n$ 在模 $m$ 意义下可逆。设 $n^{-1}$ 是它的逆元，令
\[
b\equiv an \pmod m.
\]
则当 $a$ 跑遍模 $m$ 的剩余类时，$b$ 也恰好跑遍一遍。于是
\[
\sum_{a\bmod m}\chi(a)e^{2\pi i an/m}
=\sum_{b\bmod m}\chi(n^{-1}b)e^{2\pi i b/m}.
\]
利用特征的完全乘法性，得
\[
\chi(n^{-1}b)=\chi(n^{-1})\chi(b).
\]
故上式变成
\[
\chi(n^{-1})\sum_{b\bmod m}\chi(b)e^{2\pi i b/m}
=\chi(n^{-1})\tau(\chi).
\]
而对与 $m$ 互素的整数，$|\chi(n)|=1$，故
\[
\chi(n^{-1})=\overline{\chi}(n).
\]
因此
\[
\sum_{a\bmod m}\chi(a)e^{2\pi i an/m}=\overline{\chi}(n)\tau(\chi).
\]
\end{solution}

\begin{problem}{Eisenstein与尖点分解}
\hard
设 $k\ge 4$ 为偶数，$g\in M_k(\SL_2(\Z))$ 的 $q$-展开为
\[
g(\tau)=a_0+\sum_{n\ge1}a_n q^n.
\]
证明
\[
g=a_0E_k+f,
\qquad f\in S_k(\SL_2(\Z)).
\]
并说明该分解是唯一的。
\end{problem}
\begin{solution}
因为 $E_k$ 的常数项是 $1$，所以
\[
f(\tau):=g(\tau)-a_0E_k(\tau)
\]
仍然是权 $k$ 模形式。它的常数项为
\[
a_0-a_0\cdot 1=0.
\]
因此 $f$ 在尖点 $\infty$ 处消失，也就是
\[
f\in S_k(\SL_2(\Z)).
\]
于是
\[
g=a_0E_k+f
\]
给出了所求分解。

再证唯一性。若
\[
g=cE_k+f_1=c'E_k+f_2,
\qquad f_1,f_2\in S_k,
\]
则
\[
(c-c')E_k=f_2-f_1.
\]
右边是尖点形式，常数项为 $0$；左边的常数项是 $c-c'$。所以 $c-c'=0$，即 $c=c'$，继而 $f_1=f_2$。

因此分解唯一。
\end{solution}

\begin{problem}{权十二的一组基}
\hard
利用 $M_{12}$ 是二维空间，证明
\[
M_{12}(\SL_2(\Z))=\C E_{12}\oplus \C\Delta.
\]
\end{problem}
\begin{solution}
由上一题，任意 $g\in M_{12}$ 都可唯一写成
\[
g=a_0E_{12}+f,
\qquad f\in S_{12}.
\]
因此
\[
M_{12}=\C E_{12}\oplus S_{12}.
\]
现在利用 $M_{12}$ 是二维空间，而 $E_{12}$ 已经给出一条非零向量，所以 $S_{12}$ 必是一维空间。

另一方面，$\Delta$ 是一个非零的权 $12$ 尖点形式，因此必生成整个 $S_{12}$。于是
\[
S_{12}=\C\Delta.
\]
代回即可得到
\[
M_{12}(\SL_2(\Z))=\C E_{12}\oplus \C\Delta.
\]
\end{solution}

\begin{problem}{E4三次方与E6平方的分解}
\hard
把 $E_4^3$ 与 $E_6^2$ 分别写成基
\[
\{E_{12},\Delta\}
\]
下的线性组合。
\end{problem}
\begin{solution}
因为 $E_4^3$ 与 $E_6^2$ 都属于 $M_{12}$，由上一题可写成
\[
E_4^3=E_{12}+A\Delta,
\qquad
E_6^2=E_{12}+B\Delta,
\]
这里常数项前面的系数都是 $1$，因为 $E_4^3,E_6^2,E_{12}$ 的常数项都为 $1$。

先看 $E_4^3$。它的 $q$ 一次项是 $720$，而 $E_{12}$ 的 $q$ 一次项是 $65520/691$，$\Delta$ 的首项是 $q$，故
\[
A=720-\frac{65520}{691}=\frac{432000}{691}.
\]
所以
\[
E_4^3=E_{12}+\frac{432000}{691}\Delta.
\]

再看 $E_6^2$。它的 $q$ 一次项是 $-1008$，于是
\[
B=-1008-\frac{65520}{691}=-\frac{762048}{691}.
\]
故
\[
E_6^2=E_{12}-\frac{762048}{691}\Delta.
\]

两式相减马上恢复
\[
E_4^3-E_6^2=1728\Delta.
\]
\end{solution}

\begin{problem}{Petersson正定性}
\hard
设 $f\in S_k(\SL_2(\Z))$，证明若 $f\ne 0$，则它的 Petersson 内积满足
\[
\langle f,f\rangle>0.
\]
\end{problem}
\begin{solution}
Petersson 内积定义为
\[
\langle f,g\rangle=\int_{\Gamma\backslash\HH} f(\tau)\overline{g(\tau)}\,y^k\,\frac{dx\,dy}{y^2}.
\]
取 $g=f$，便有
\[
\langle f,f\rangle=\int_{\Gamma\backslash\HH}|f(\tau)|^2y^k\,\frac{dx\,dy}{y^2}.
\]
被积函数处处非负。

由于 $f$ 是尖点形式，它在尖点处指数衰减，所以该积分收敛。若 $f\ne 0$，则 $|f(\tau)|^2$ 不是处处为零，因此在某个开集上严格为正。乘上正函数 $y^k$ 后仍严格为正，所以积分值严格大于零。

因此对任意非零尖点形式都有
\[
\langle f,f\rangle>0.
\]
\end{solution}

\begin{problem}{带特征的展开}
\hard
设 $\chi_4$ 为模 $4$ 的特征。定义系数
\[
a_n=\sum_{d\mid n}\chi_4(d)d^2.
\]
写出级数
\[
F(\tau)=\sum_{n\ge1} a_n q^n
\]
的前五项，并说明为什么这里选权 $3$ 是自然的。
\end{problem}
\begin{solution}
先逐项计算系数。

当 $n=1$ 时，因子只有 $1$，故
\[
a_1=\chi_4(1)\cdot 1^2=1.
\]

当 $n=2$ 时，因子为 $1,2$，而 $\chi_4(2)=0$，故
\[
a_2=1.
\]

当 $n=3$ 时，因子为 $1,3$，于是
\[
a_3=\chi_4(1)\cdot 1^2+\chi_4(3)\cdot 3^2=1-9=-8.
\]

当 $n=4$ 时，只有奇因子 $1$ 有贡献，所以
\[
a_4=1.
\]

当 $n=5$ 时，因子为 $1,5$，且 $\chi_4(5)=\chi_4(1)=1$，故
\[
a_5=1+25=26.
\]
因此
\[
F(\tau)=q+q^2-8q^3+q^4+26q^5+\cdots.
\]

为什么权 $3$ 自然？因为模 $4$ 的特征满足
\[
\chi_4(-1)=-1.
\]
而带特征 Eisenstein 级数的奇偶性要求权 $k$ 与特征满足
\[
\chi(-1)=(-1)^k.
\]
因此这里应选奇数权；最小的正整数选择就是 $k=3$。
\end{solution}

\begin{problem}{Ramanujan系数的卷积公式}
\hard
设
\[
\Delta(\tau)=\sum_{n\ge1}\tau(n)q^n,
\qquad
A(\tau)=\sum_{n\ge1}\sigma_3(n)q^n,
\qquad
B(\tau)=\sum_{n\ge1}\sigma_5(n)q^n.
\]
由
\[
E_4=1+240A,
\qquad
E_6=1-504B,
\qquad
E_4^3-E_6^2=1728\Delta
\]
导出 $\tau(n)$ 的一个显式卷积公式。
\end{problem}
\begin{solution}
先把两边完全展开。由
\[
E_4^3=(1+240A)^3=1+720A+172800A^2+13824000A^3,
\]
以及
\[
E_6^2=(1-504B)^2=1-1008B+254016B^2,
\]
可得
\[
1728\Delta
=720A+1008B+172800A^2+13824000A^3-254016B^2.
\]
取 $q^n$ 的系数。记
\[
(A^2)_n=\sum_{a+b=n}\sigma_3(a)\sigma_3(b),
\]
\[
(A^3)_n=\sum_{a+b+c=n}\sigma_3(a)\sigma_3(b)\sigma_3(c),
\]
\[
(B^2)_n=\sum_{a+b=n}\sigma_5(a)\sigma_5(b),
\]
其中求和都在正整数上进行。于是得到
\[
1728\,\tau(n)
=720\sigma_3(n)+1008\sigma_5(n)+172800(A^2)_n+13824000(A^3)_n-254016(B^2)_n.
\]
也就是说
\[
\tau(n)=\frac{1}{1728}\Bigl(
720\sigma_3(n)+1008\sigma_5(n)
+172800\sum_{a+b=n}\sigma_3(a)\sigma_3(b)
\Bigr.
\]
\[
\Bigl.
+13824000\sum_{a+b+c=n}\sigma_3(a)\sigma_3(b)\sigma_3(c)
-254016\sum_{a+b=n}\sigma_5(a)\sigma_5(b)
\Bigr).
\]
这就是由 Eisenstein 级数恒等式得到的 Ramanujan 系数卷积公式。
\end{solution}
